\section*{Розділ IV. Керівні органи ЦВКс}
\addcontentsline{toc}{section}{Розділ IV. Керівні органи ЦВКс}

\subsection*{4.1. Голова ЦВКс}
\addcontentsline{toc}{subsection}{4.1. Голова ЦВКс}
    4.1.1. ЦВКс очолює Голова ЦВКс, який обирається членами ЦВКс зі свого складу простою більшістю голосів на першому засіданні після формування комісії.

    4.1.2. До повноважень Голови ЦВКс належить:

        \begin{enumerate}[label=\alph*)]
            \item Скликання засідань ЦВКс та визначення їх порядку денного;
            \item Головування на засіданнях ЦВКс;
            \item Організація виконання рішень ЦВКс;
            \item Представлення ЦВКс у відносинах з органами студентського самоврядування, адміністрацією Університету та іншими суб'єктами;
            \item Підписання рішень, протоколів та інших документів ЦВКс;
            \item Координація роботи членів ЦВКс під час виборчих кампаній;
            \item Прийняття оперативних рішень між засіданнями ЦВКс з подальшим інформуванням комісії;
            \item Забезпечення дотримання строків проведення виборчих процедур.
        \end{enumerate}

    4.1.3. Голова ЦВКс має право:

        \begin{enumerate}[label=\alph*)]
            \item Скликати позачергові засідання ЦВКс;
            \item Ставити питання на голосування;
            \item Запрошувати на засідання ЦВКс представників органів студентського самоврядування, адміністрації Університету та експертів;
            \item Призупиняти виборчий процес у разі виявлення серйозних порушень до розгляду питання на засіданні ЦВКс.
        \end{enumerate}

    4.1.4. Голова ЦВКс зобов'язаний:

        \begin{enumerate}[label=\alph*)]
            \item Забезпечувати колегіальність у прийнятті рішень;
            \item Дотримуватися принципів неупередженості та справедливості;
            \item Інформувати членів ЦВКс про всі оперативні рішення, прийняті між засіданнями;
            \item Подавати звіти про діяльність ЦВКс до КСУ.
        \end{enumerate}

\subsection*{4.2. Секретар ЦВКс}
\addcontentsline{toc}{subsection}{4.2. Секретар ЦВКс}
    4.2.1. Секретар ЦВКс обирається членами ЦВКс зі свого складу простою більшістю голосів на першому засіданні після формування комісії.

    4.2.2. До повноважень Секретаря ЦВКс належить:

        \begin{enumerate}[label=\alph*)]
            \item Ведення протоколів засідань ЦВКс;
            \item Забезпечення документообігу ЦВКс;
            \item Ведення архіву документів ЦВКс;
            \item Підготовка засідань ЦВКс (розсилка повідомлень, підготовка матеріалів);
            \item Координація роботи з виборчими дільницями та комісіями нижчого рівня;
            \item Забезпечення технічної організації виборчих процедур;
            \item Ведення обліку кандидатів та виборців;
            \item Підготовка звітів про діяльність ЦВКс.
        \end{enumerate}

    4.2.3. Секретар ЦВКс має право:

        \begin{enumerate}[label=\alph*)]
            \item Запитувати необхідну інформацію від органів студентського самоврядування для виконання своїх обов'язків;
            \item Брати участь у підготовці рішень ЦВКс;
            \item Координувати роботу технічних працівників, залучених до виборчого процесу.
        \end{enumerate}

\subsection*{4.3. Заміщення Голови ЦВКс}
\addcontentsline{toc}{subsection}{4.3. Заміщення Голови ЦВКс}
    4.3.1. ЦВКс на своєму першому засіданні обирає заступника Голови ЦВКс зі складу членів ЦВКс.

    4.3.2. Заступник Голови ЦВКс виконує функції Голови у разі:

        \begin{enumerate}[label=\alph*)]
            \item Відсутності Голови з поважних причин;
            \item Неможливості Голови виконувати свої обов'язки;
            \item Дострокового припинення повноважень Голови до обрання нового.
        \end{enumerate}

    4.3.3. Заступник Голови ЦВКс має всі повноваження Голови під час виконання його функцій та несе повну відповідальність за прийняті рішення.

\subsection*{4.4. Переобрання керівництва ЦВКс}
\addcontentsline{toc}{subsection}{4.4. Переобрання керівництва ЦВКс}
    4.4.1. Голова, Секретар або заступник Голови ЦВКс можуть бути переобрані на нових виборах за рішенням більшості членів ЦВКс у разі:

        \begin{enumerate}[label=\alph*)]
            \item Незадовільного виконання своїх обов'язків;
            \item Систематичного порушення вимог цього Положення;
            \item Втрати довіри більшості членів ЦВКс;
            \item Подання особистої заяви про складення повноважень.
        \end{enumerate}

    4.4.2. Рішення про переобрання керівництва ЦВКс приймається таємним голосуванням за ініціативою не менше половини членів ЦВКс.

    4.4.3. Нове обрання на посаду проводиться на тому ж засіданні ЦВКс, на якому прийнято рішення про переобрання.

\subsection*{4.5. Кворум та процедури прийняття рішень}
\addcontentsline{toc}{subsection}{4.5. Кворум та процедури прийняття рішень}
    4.5.1. Засідання ЦВКс вважається правомочним за наявності простої більшості від загального складу членів ЦВКс.

    4.5.2. Рішення ЦВКс приймаються відкритим голосуванням, крім випадків, коли ЦВКс приймає рішення про таємне голосування.

    4.5.3. Процедура прийняття рішень ЦВКс диференціюється залежно від характеру питання:

        \begin{enumerate}[label=\alph*)]
            \item Поточні організаційні рішення (призначення засідань, розподіл обов'язків, технічні питання) приймаються простою більшістю голосів від присутніх членів ЦВКс;
            \item Принципові рішення (скасування реєстрації кандидата, визнання виборів недійсними, зміна виборчих процедур, припинення виборчого процесу) приймаються більшістю у дві третини (2/3) голосів від загального складу ЦВКс.
        \end{enumerate}

    4.5.4. У разі рівного розподілу голосів при прийнятті поточних рішень, вирішальним є голос Голови ЦВКс.

\subsection*{4.6. Відповідальність керівних органів}
\addcontentsline{toc}{subsection}{4.6. Відповідальність керівних органів}
    4.6.1. Голова, заступник Голови та Секретар ЦВКс несуть персональну відповідальність за:

        \begin{enumerate}[label=\alph*)]
            \item Своєчасне та якісне виконання своїх функціональних обов'язків;
            \item Дотримання принципів неупередженості та справедливості;
            \item Збереження конфіденційної інформації, отриманої під час виконання обов'язків;
            \item Забезпечення прозорості та законності виборчого процесу.
        \end{enumerate}

    4.6.2. Колегіальну відповідальність за діяльність ЦВКс несуть усі її члени.